% Section 3: Forecasting for reinforcement learning
\subsection{Forecasting for reinforcement learning}
\label{sec:forecasting_for_rl}

\subsubsection{World models}
\label{sec:world_models}

The dominant paradigm for using forecasting inside reinforcement learning is the
world model. The architecture in its modern form is due to \citet{HaSchmidhuber2018},
who decompose the agent into three components, a vision module $V$ that compresses
the raw observation $x_t$ into a low-dimensional latent code $z_t \in \mathbb R^{32}$
via a variational autoencoder, a memory module $M$ that predicts the next latent
through a mixture density network on top of a recurrent neural network, and a
controller module $C$ that maps $(z_t, h_t)$ to the next action. The memory model
parametrises the conditional distribution
$P(z_{t+1} \mid a_t, z_t, h_t)$ as a temperature-scaled mixture of Gaussians, where
$h_t$ is the hidden state of the recurrent network and the temperature $\tau$ acts
as a stochasticity control during sampling. The controller is the linear map
$a_t = W_c [z_t; h_t] + b_c$, deliberately under-parametrised so that the model
complexity sits in $V$ and $M$ and is trained by Covariance Matrix Adaptation
Evolution Strategies because the loss surface is not differentiable through the
environment. Parameter counts on CarRacing-v0 are $4.35$ million for the VAE,
$0.42$ million for the MDN-RNN, and $867$ for the controller.

The two empirical demonstrations of \citet{HaSchmidhuber2018} are CarRacing-v0 and
the VizDoom Take Cover environment. On CarRacing-v0, the full world-model agent
attains an average score of $906 \pm 21$ over $100$ trials, which is the first
reported solution to the task at the OpenAI Gym threshold of $900$ and exceeds the
prior state of the art at $838 \pm 11$. On VizDoom Take Cover, the more striking
result, the world model is trained on $10{,}000$ rollouts of a random policy, and
the controller is then trained \emph{entirely inside the learned dream}, namely a
hallucinated environment generated by $M$ without any access to the real game
engine. The controller transferred zero-shot to the real environment achieves an
average survival score of approximately $1{,}100$ time steps at temperature
$\tau = 1.15$, against a solve threshold of $750$ and the Gym leaderboard at
$820 \pm 58$. The dream-trained controller actually outperforms the leaderboard.
The technical caveat is the exploitation problem, namely that the controller may
discover adversarial policies that look optimal under $M$ but fail in reality.
\citet{HaSchmidhuber2018} mitigate this by training $M$ as a stochastic mixture
density network and by tuning $\tau$ upward to make the dream environment harder.

The Dreamer family of \citet{Hafner2025} refines the architecture into the
Recurrent State Space Model and scales it to a single agent that succeeds on more
than $150$ tasks across eight diverse domains. The world model combines a
deterministic hidden state $h_t$ and a stochastic categorical representation $z_t$
sampled with straight-through gradients. Three losses are minimised end-to-end
with fixed weights $\beta_{\mathrm{pred}} = 1$, $\beta_{\mathrm{dyn}} = 1$,
$\beta_{\mathrm{rep}} = 0.1$,
\begin{equation}
\mathcal L_{\mathrm{world}}(\phi)
= \mathbb E\Bigl[\sum_t
  \beta_{\mathrm{pred}}\, \mathcal L_{\mathrm{pred}}
  + \beta_{\mathrm{dyn}}\, \mathcal L_{\mathrm{dyn}}
  + \beta_{\mathrm{rep}}\, \mathcal L_{\mathrm{rep}}
\Bigr],
\label{eq:dreamer_world_loss}
\end{equation}
where the prediction loss covers image reconstruction, reward prediction via a
symlog squared loss, and the binary continuation flag via logistic regression, the
dynamics loss is the Kullback-Leibler divergence between the prior
$p_\phi(z_t \mid h_t)$ and the stop-gradient posterior $q_\phi(z_t \mid h_t, x_t)$,
and the representation loss is the symmetric KL term. Free bits clip the dynamics
and representation losses below one nat to avoid degenerate solutions, and the
categorical distributions are parametrised as $99\%$ neural network output mixed
with $1\%$ uniform to keep the KL terms well behaved.

The actor and critic operate on imagined trajectories of length $T = 16$ generated
by the world model. The critic learns a categorical distribution over
exponentially-spaced return bins to handle reward scales that vary by orders of
magnitude across domains, the actor uses the REINFORCE estimator with a fixed
entropy regulariser $\eta = 3 \times 10^{-4}$, and the discount factor is
$\gamma = 0.997$. Returns are normalised by the range from the $5$th to the
$95$th percentile so that the entropy regulariser carries a comparable weight
under both dense and sparse rewards. The symlog transformation $f(x) =
\mathrm{sign}(x)\, \log(|x| + 1)$, applied to vector observations and to the
reward predictor's targets, is the single most consequential engineering choice
in \citet{Hafner2025}, because it compresses target magnitudes symmetrically
around zero without truncating large values and without introducing
non-stationarity from running normalisation.

The empirical scope of \citet{Hafner2025} is the broadest in the world-model
literature. On Atari (57 games, 200M frames) Dreamer outperforms MuZero, Rainbow,
and IQN at a fraction of the compute. On Atari100k (26 games, 400K frames)
Dreamer beats the transformer-based IRIS and TWM, the model-free SPR, and SimPLe.
On DMLab (30 tasks, 100M frames) Dreamer matches IMPALA and R2D2 at 1B steps,
a tenfold improvement in data efficiency. On Proprio Control (18 tasks, 500K
steps) and Visual Control (20 tasks, 1M steps) Dreamer is the new state of the
art across D4PG, DMPO, MPO, DrQ-v2, and CURL. On Minecraft, the long-standing
diamond-collection challenge, Dreamer is the first algorithm to collect diamonds
from scratch in $100$ million environment steps without human expert data and
without an adaptive curriculum. The ablation in Figure~6 of \citet{Hafner2025}
attributes most of the gain to the symlog and twohot prediction tools, with the
representation-loss KL objective second, and reports the surprising finding that
the agent's performance rests predominantly on the unsupervised reconstruction
loss rather than on reward and value prediction gradients.

The methodological message for our chapter is that a learned probabilistic
forecaster of dynamics can replace a hand-built simulator across a remarkably
wide range of domains, and that the engineering tricks required to make this work
out of the box, namely robust target transformations, free bits, return
percentile normalisation, and the actor-critic-on-imagination loop, are now
well enough understood to ship as a single hyperparameter configuration. The
implication for the economics audience is that the world-model recipe is
mature on the dynamics side, and the open question is whether a world model can
absorb an action-conditional time series foundation model on the forecasting
side without losing the calibration properties on which decision-aware training
depends. We return to this question in Section~\ref{sec:tsfm} and
Section~\ref{sec:forecasting_rl_open}.

\subsubsection{Trajectory generative models}
\label{sec:trajectory_models}

A parallel line of work treats the trajectory itself as the object to be modelled,
and replaces the value-function and policy-gradient apparatus of conventional
offline reinforcement learning with the apparatus of large-scale sequence
generation. The Decision Transformer of \citet{Chen2021DT} recasts offline RL as
return-conditional autoregressive sequence modelling. Trajectories are represented
as a sequence of triples
$\tau = (\hat R_1, s_1, a_1, \hat R_2, s_2, a_2, \dots, \hat R_T, s_T, a_T)$, where
$\hat R_t = \sum_{t' \ge t} r_{t'}$ is the return-to-go from time $t$ onward and
serves as the conditioning prompt at inference. A causally masked GPT-style
transformer with linear modality embeddings and per-timestep positional
embeddings predicts the next action given the last $K$ context tokens. At test
time, the operator specifies a desired return $R^\star$, primes the model with
$\hat R_1 = R^\star$ and the current state, samples the action autoregressively,
executes it, decrements $R^\star$ by the realised reward, and repeats. No
dynamic programming, no policy gradient, and no critic appear at training time.

The empirical case for the Decision Transformer rests on direct comparisons with
Conservative Q-Learning \citep{Kumar2020} on offline benchmarks. On the 1\% DQN
replay slice of Atari, the Decision Transformer attains gamer-normalised scores of
$267.5$ on Breakout, $106.1$ on Pong, and $2.5$ on Seaquest, against CQL's
$211.1$, $111.9$, and $1.7$ on the same games, with the Decision Transformer
losing only on Qbert. On the D4RL locomotion benchmark, the average normalised
return across the medium, medium-replay, and medium-expert subsets is $74.7$ for
the Decision Transformer against $54.2$ for CQL, $48.2$ for BEAR, $36.9$ for
BRAC-v, $34.3$ for AWR, and $47.7$ for behavioural cloning. The diagnostic
analysis of \citet{Chen2021DT} establishes that the Decision Transformer is not
performing behavioural cloning on a top-percentile subset of the data, that
longer context windows help, and that the architecture remains effective in
sparse-reward Key-to-Door tasks where conventional credit assignment via Bellman
backups fails.

The Trajectory Transformer of \citet{Janner2021TT} replaces return conditioning
with discrete tokenisation and beam-search planning. The trajectory is flattened
into a sequence of length $T(N + M + 2)$ where each state and action dimension is
discretised independently into either uniform-width or quantile bins, and reward
and return-to-go follow each transition. A four-layer four-head GPT-style decoder
is trained by teacher forcing on the autoregressive log-likelihood. At test time,
beam search with width $B$ replaces transition log-probabilities with predicted
reward plus return-to-go, which converts the trajectory sequence model into a
reward-maximising planner that does not exploit dynamics out of distribution.
On D4RL locomotion, the quantile-discretisation variant attains an average of
$78.9$ against $74.7$ for the Decision Transformer and $77.6$ for CQL. The more
striking result is on the AntMaze sparse-reward navigation suite. Augmenting the
Trajectory Transformer's beam search with an implicit-Q-learning critic as the
value heuristic, the combined planner attains an average normalised score of
$84.0$ across the five maze sizes against $63.2$ for IQL alone, $44.9$ for CQL,
and $11.8$ for the Decision Transformer. The Q-guided planner outperforms IQL
itself, which is the policy from which the Q-function was extracted, an unusual
result that \citet{Janner2021TT} attribute to the temporal compositionality of
beam search over dynamic-programming targets.

Diffuser \citep{Janner2022diffuser} pushes the abstraction further by replacing
autoregressive generation with iterative denoising of an entire trajectory. The
trajectory is represented as a two-dimensional array
$\tau \in \mathbb R^{(d_s + d_a) \times T}$ with one column per planning step, the
forward process is a standard Gaussian noise schedule, and the reverse process is
parametrised by a temporal U-Net whose receptive field is restricted to local
windows in the planning-time dimension. Local consistency of denoising is composed
into global coherence over many denoising steps. The reinforcement learning problem
is recast as inference over a perturbed trajectory distribution
\begin{equation}
\tilde p(\tau) \propto p_\theta(\tau)\, h(\tau),
\label{eq:diffuser_perturbed}
\end{equation}
where $h(\tau)$ is a return-related guide that biases the sampler toward
high-value trajectories. The classifier-guided variant
$\tau^{i-1} \sim \mathcal N\bigl( \mu_\theta(\tau^i) + \alpha\, \Sigma^i\,
\nabla J_\phi(\mu_\theta(\tau^i)),\, \Sigma^i \bigr)$
inserts the gradient of a learned return predictor $J_\phi$ into the standard
reverse-process mean, while the inpainting variant pins the initial state and any
goal state by replacing the corresponding columns at each denoising step. The
fully convolutional horizon dimension lets the planning length be a runtime
parameter, the joint state-action representation prevents the dynamics from being
queried at adversarial action sequences, and the same Diffuser model serves
multiple reward functions via different guide functions $h$. On the long-horizon
Maze2D benchmark, where shooting-based planners typically fail, Diffuser
substantially outperforms model-free baselines, and on D4RL locomotion the gap
with the Trajectory Transformer narrows but remains in favour of the diffusion
approach on the medium-replay subset.

The Decision Diffuser of \citet{Ajay2023} pushes the line further still by
making the policy itself a return-conditional diffusion model and abandoning
both the value-function machinery and the joint state-action representation. The
diffusion is taken over states alone, namely
$x^k(\tau) = (s_1, \dots, s_H)^k$ in the diffusion ladder, and the action is
extracted from consecutive predicted states via an inverse-dynamics model
$a_t = f_\phi(s_t, s_{t+1})$ trained on the same offline data. Conditioning uses
classifier-free guidance, namely the perturbed noise prediction
\begin{equation}
\hat\varepsilon = \varepsilon_\theta(x^k, k)
  + \omega\,\bigl(\varepsilon_\theta(x^k, y, k)
    - \varepsilon_\theta(x^k, \emptyset, k)\bigr),
\label{eq:dd_classifier_free}
\end{equation}
with $\omega$ the guidance scale and $y$ either the normalised return, a
constraint indicator, or a skill indicator. The classifier-free formulation
avoids estimating a Q-function and therefore avoids the implicit dynamic
programming that the rest of the offline-RL toolkit relies on. The empirical
case is that the Decision Diffuser is competitive with or beats CQL, IQL, the
Trajectory Transformer, the Decision Transformer, and the original Diffuser
across D4RL locomotion, that the gap widens on the long-horizon D4RL Kitchen
tasks, and that the same model satisfies novel combinations of constraints in
the Kuka block-stacking environment and composes quadruped gaits in the Unitree
environment after training only on single constraints or single skills. The
methodological message across this strand is that offline reinforcement learning
can be reduced to conditional generation of high-dimensional structured objects,
with the forecast distribution over trajectories serving simultaneously as the
substrate for policy evaluation, policy improvement, and constraint satisfaction.

\subsubsection{Counterfactual forecasting for off-policy evaluation}
\label{sec:counterfactual_ope}

Off-policy evaluation in finite samples is bias-variance constrained. Importance
sampling estimators are unbiased under common support but high variance, while
model-based estimators are low variance but biased under model misspecification.
A third route, introduced by \citet{Buesing2019}, recasts the partially-observed
Markov decision process as a structural causal model and uses counterfactual
inference rather than interventional simulation as the off-policy evaluation
estimator. A structural causal model over variables $X = (X_1, \dots, X_N)$ is
a triple $(\mathcal G, \{f_i\}, \{P_{U_i}\})$ consisting of a directed acyclic
graph, structural mechanisms, and independent noise distributions, such that
$X_i = f_i(\mathrm{pa}_i, U_i)$. A POMDP fits the framework with initial state
$U^s_1 = S_1$, transition mechanism $S_t = f^s(S_{t-1}, A_{t-1}, U^s_t)$,
observation mechanism $O_t = f^o(S_t, U^o_t)$, and action mechanism
$A_t = f^\pi(H_t, U^a_t)$ given the history $H_t$.

The counterfactual inference rule of \citet{Buesing2019} consists of three steps,
namely abduction, action, and prediction. Given an observed trajectory
$\hat\tau$, infer the posterior $P(U \mid \hat\tau)$ over the exogenous noise.
Then replace the prior $P_U$ with the posterior $P(U \mid \hat\tau)$ in the SCM,
denoted $\mathcal M_{\hat\tau}$. Finally, apply the intervention
$\mathrm{do}(I)$, namely substituting the behaviour-policy mechanism with the
target policy $\pi$, and read off the implied counterfactual distribution
$P^{\mathrm{do}(I) \mid \hat\tau}$. Lemma~1 of \citet{Buesing2019} states that
marginalising the counterfactual density over observations recovers the
interventional density, namely
$\mathbb E_{\hat\tau \sim P}\bigl[ P^{\mathrm{do}(I) \mid \hat\tau}(x_q) \bigr]
= P^{\mathrm{do}(I)}(x_q)$,
which by Corollary~2 implies that the counterfactual policy-evaluation estimator
is unbiased under correct model specification. The empirical case for
counterfactual policy evaluation over model-based policy evaluation is that
$P(U \mid \hat\tau)$ concentrates on the parts of noise space consistent with
real episodes, so model errors that propagate from $P_U$ to the evaluation
estimate under interventional simulation do not propagate under counterfactual
simulation when the dynamics mechanism is accurate but the prior on noise is
not. On the PO-SOKOBAN grid-world environment of \citet{Buesing2019}, the
counterfactual estimator's evaluation error decreases monotonically with the
amount of off-policy data, while the model-based estimator's error remains
stuck at the bias of the model.

The categorical-action case raises an identification problem that
\citet{Buesing2019} sidestep by working in a deterministic-transition grid world
and that \citet{Oberst2019} address head on. The non-identifiability is that
multiple structural causal models can be consistent with the same interventional
distribution yet imply different counterfactual outcomes. The example of
\citet{Oberst2019} considers a categorical transition with $k = 4$ outcomes and
two alternative orderings of the unit interval that both sample from the
prescribed distribution, but that imply different counterfactual states under
the same posterior. The Pearl monotonicity condition resolves the
non-identifiability for binary outcomes, and \citet{Oberst2019} introduce the
analogous condition for categorical variables, namely counterfactual stability,
which requires that if outcome $i$ was observed under intervention $I$ then the
counterfactual outcome under $I'$ can be $j \ne i$ only if the multiplicative
change in $p_j$ strictly exceeds the multiplicative change in $p_i$.

Their Theorem~1 shows that counterfactual stability for $k = 2$ reduces to
Pearl's monotonicity. Their constructive proposal is the Gumbel-max structural
causal model, which samples a categorical outcome via
$Y = \arg\max_j \{ \log p_j + G_j \}$ for independent Gumbel noise variables
$G_j$, and they prove in Theorem~2 that this SCM satisfies counterfactual
stability. Posterior inference under the Gumbel-max model is then explicit. The
observed argmax $i$ pins down a truncated joint of the Gumbel noise, and
counterfactual sampling proceeds by drawing fresh Gumbel variables consistent
with the truncated joint and re-evaluating the argmax under the alternative
probabilities. The application is a sepsis management POMDP, where the
Gumbel-max counterfactual rule highlights specific episodes on which the
candidate RL policy would have diverged from the physician policy by the largest
margin, which is the right diagnostic for a clinician auditing a proposed
treatment rule.

\citet{TangWiens2023} bridge the gap between the pure counterfactual route of
\citet{Buesing2019} and \citet{Oberst2019} and the importance-sampling
tradition. They allow the dataset to be augmented with human-provided
counterfactual annotations $g^{\tilde a}_t \sim G_t(s_t, \tilde a)$ that estimate
the future return $Q^{\pi_e}_{t:T}(s_t, \tilde a)$ under the evaluation policy
$\pi_e$ for an alternative action $\tilde a$, and they construct the
counterfactual-augmented per-decision importance sampling estimator
\begin{equation}
\hat v_{\mathrm{C\text{-}PDIS}}
= v_T,
\quad
v_{T-t+1}
= w^{a_t}_t \rho^{a_t}_t (r_t + \gamma v_{T-t})
+ \sum_{\tilde a \ne a_t} w^{\tilde a}_t \rho^{\tilde a}_t g^{\tilde a}_t,
\label{eq:tangwiens_cpdis}
\end{equation}
with weights $w^{\tilde a}_t$ summing to one across actions at each step and
importance ratios $\rho^{\tilde a}_t = \pi_e(\tilde a \mid s_t) /
\pi_{b+}(\tilde a \mid s_t)$ relative to an augmented behaviour policy that
includes the annotations. Theorem~1 establishes that under perfect annotation
and an extended common-support condition, C-IS is unbiased in the bandit case,
and Theorem~8 extends the result to MDPs by backward induction. Theorem~6
gives a closed-form variance decomposition under equal annotation weights,
and Proposition~7 establishes a guaranteed variance reduction over standard
importance sampling. On a sepsis simulator with $50$ replications of a
$1{,}000$-episode dataset, the C-PDIS estimator with equal weights and perfect
$Q^{\pi_e}$ annotations attains an out-of-policy RMSE of $0.013$ against $0.113$
for vanilla PDIS, an effective sample size of $994.0$ out of $1{,}000$ against
$76.8$, and a Spearman policy-ranking correlation of $0.995$ against $0.596$.
Under the more realistic regime where the annotations summarise the behaviour
policy's $Q^{\pi_b}$ rather than the target's $Q^{\pi_e}$, the RMSE rises to
$0.070$ and the Spearman correlation falls to $0.961$, both still well above
the PDIS baseline.

The shared identifying assumption across this strand is that the underlying
structural causal model is correctly specified, namely that the postulated
exogenous noise is independent of past observations conditional on the state
and action and that the categorical-counterfactual stability condition holds.
The Lucas critique applies here in its sharpest form. If the agent's policy
itself enters the data-generating process, namely if the SCM's noise marginals
$P_U$ depend on the deployed policy through equilibrium considerations, then
counterfactual estimates of an alternative policy's value will be biased by
exactly the structural shift that economic agents respond to. The credible
remedies, namely instrumental-variable identification and exogeneity
restrictions imported from the machine-learning-instruments literature, are
the subject of Section~\ref{sec:forecasting_rl_open}.

\subsubsection{Time series foundation models}
\label{sec:tsfm}

The most recent development on the forecasting-for-RL side is the arrival of
pretrained univariate time series foundation models. The three architectural
templates in current use differ in how they represent the time series for a
transformer backbone, namely categorical tokenisation, continuous patching, and
explicit lag covariates. Each template has advanced the state of zero-shot
forecasting accuracy on benchmark suites, and the natural question for our
purposes is whether any of them can serve as the action-conditional dynamics
model an RL agent requires. The answer is no, for reasons we lay out at the end
of this subsection. The contributions of the four supporting papers in this
strand are most usefully read against that conclusion.

Chronos of \citet{Ansari2024} treats time series forecasting as a language
modelling problem. Given a univariate series $x_{1:C+H}$, the historical context
$x_{1:C}$ is first mean-scaled by $s = C^{-1} \sum_i |x_i|$ and then quantised
into a fixed vocabulary of $B$ uniformly-spaced bins. The quantisation map
$q : \mathbb R \to \{1, \dots, B\}$ assigns each scaled value to its enclosing
bin, and the dequantisation map $d : \{1, \dots, B\} \to \mathbb R$ returns the
bin centre. The tokenised series, together with PAD and EOS markers, is fed to
a T5 encoder-decoder or GPT-2 decoder-only transformer with the vocabulary
truncated to size $|V_{\mathrm{ts}}|$, and the model is trained with the standard
language-modelling cross-entropy
\begin{equation}
\mathcal L(\theta) = - \sum_{h=1}^{H+1} \sum_{i=1}^{|V_{\mathrm{ts}}|}
  \mathbf 1\{z_{C+h+1} = i\}\, \log p_\theta(z_{C+h+1} = i \mid z_{1:C+h}).
\label{eq:chronos_loss}
\end{equation}
The Chronos family covers T5 backbones from $20$M to $710$M parameters with
$B = 4{,}096$ at the larger end, and the training corpus is augmented by
TSMixup, namely convex combinations of randomly sampled series, and by
KernelSynth, namely Gaussian-process synthetic series with composed
linear-plus-RBF-plus-periodic kernels. The forecast at inference time is a
sample-based predictive distribution over the next bin, mapped back to real
values, which gives Chronos native support for probabilistic forecasting.

TimesFM of \citet{Das2024} works in the continuous domain and treats patches of
consecutive time points rather than individual values as the unit of
representation. The input series of length $L$ is divided into non-overlapping
patches of length $p$, each patch is processed by a residual block into a vector
of size $d_{\mathrm{model}}$, and the resulting $N = \lfloor L / p \rfloor$
tokens are fed to a decoder-only causal-attention transformer of $200$M
parameters. The output of patch $j$ is a forecast of the next
$\mathrm{output\_patch\_len}$ time points, so the output horizon need not equal
the input patch length and long-horizon forecasts can be generated without
multi-step autoregressive decoding. The training objective is mean squared error
on the realised future patches, and the pretraining corpus consists of
approximately $100$ billion time-point observations drawn from Google Trends,
Wiki Pageviews, and synthetic ARMA, seasonal, trend, and step-function
processes. \citet{Das2024} show that out-of-the-box zero-shot performance on a
diverse set of unseen forecasting datasets comes close to the accuracy of
fully supervised models trained for each individual dataset.

Lag-Llama of \citet{Rasul2023} eschews both tokenisation and patching in favour
of an explicit lag-feature representation that maps directly onto the classical
autoregressive structure familiar to time-series econometrics. At each timestep
$t$, the input vector concatenates the lag values
$\{x_{t-\ell} : \ell \in \mathcal L\}$ for a sorted set of lag indices spanning
quarterly, monthly, weekly, daily, hourly, and second-level frequencies,
together with date-time covariates and per-window summary statistics
$\mu_i, \sigma_i$. The backbone is a decoder-only transformer with RMSNorm
pre-normalisation and rotary positional encodings, and the output head is a
parametric distribution, typically a Student's-$t$, optimised by negative log
likelihood. Robust standardisation, namely subtracting the window median and
dividing by the interquartile range, replaces the standard mean-variance
scaling. Pretraining uses a corpus of $7{,}965$ series across $27$ datasets and
six semantically-grouped domains, with approximately $352$ million data windows
in total. Lag-Llama's appeal to economists is structural. The lag representation
is what an econometrician would write down as the natural input for a
generalised autoregression, and the parametric output head is a direct analogue
of the Student's-$t$ error specification used in financial time-series models.

The form-and-function critique of \citet{PerezDiaz2025} identifies the central
operational limitation of all three architectures. They formalise the four
forecast types produced by current TSFMs, namely point forecasts, quantile
forecasts, per-step parametric marginals, and trajectory ensembles, and prove
that the expressiveness hierarchy is strict. Trajectory ensembles can be
marginalised to any of the other three, but the reverse direction requires
imposing a copula or another dependence structure. They then identify six
canonical operational tasks, namely pointwise prediction intervals, pathwise
forecast bands, event and tail probabilities, threshold crossing and persistence,
window aggregates, and scenario generation, and establish a compatibility matrix
that shows which forecast types natively support each task. Path-dependent
queries, including threshold crossings and window aggregates, are unsupported by
point, quantile, and parametric forecasts without additional dependence
assumptions. The complementary scoring-rule analysis distinguishes marginally
proper rules such as the continuous ranked probability score from jointly proper
rules such as the energy score, and shows that low CRPS does not imply low
energy score for the same model. The implication for an RL agent that needs to
reason about path-dependent events, including threshold-crossing inventory
decisions and worst-case execution costs, is that the marginal forecasts
produced by most current TSFMs are insufficient, and the only forecast type
that directly supports such queries is the trajectory-ensemble output, which
remains less common in the TSFM literature.

The negative result of \citet{Rahimikia2025} provides the strongest existing
empirical counterpoint to the zero-shot-foundation-model narrative. They
evaluate Chronos and TimesFM on a panel of daily excess returns spanning $34$
years across $94$ countries with approximately two billion observations, and
compare against linear, Lasso, Ridge, neural network, and tree-ensemble
baselines. The headline numbers are stark. Off-the-shelf zero-shot Chronos
(large) with a $512$-step context attains an out-of-sample $R^2$ of $-1.37\%$
and a directional accuracy just above $51\%$. Zero-shot TimesFM ($500$M) attains
$R^2 = -2.80\%$ and directional accuracy just below $50\%$. The annualised
long-short portfolio returns at a $512$ context are $20.17\%$ for Chronos
(large) and $-1.47\%$ for TimesFM ($500$M). Fine-tuning the pretrained models on
financial data closes the $R^2$ gap only partially and does not always translate
into economic gains. Pretraining from scratch on financial data, however, does
realise the full economic potential of the architecture and outperforms the
tree-ensemble baselines on the goodness-of-fit metrics. The lesson is that
generic time-series pretraining does not transfer to financial domains, and
that finance-native pretraining and data scaling are essential. For an
RL-in-finance application, the implication is that an off-the-shelf TSFM is not
ready to serve as a world model for return processes.

\citet{Jetwiriyanon2025} examine a complementary question on the slower-moving
macroeconomic side, namely whether the headline TSFM accuracy claims survive in
the noisy, low-sample-size regime that characterises GDP forecasting. They
evaluate Chronos, Moirai, and TimeGPT on a quarterly New Zealand GDP series
spanning 1999Q3 to 2024Q3, against the Reserve Bank of New Zealand's LSBoost and
factor-model benchmarks and against persistence and ARIMA baselines. The
empirical finding is that zero-shot TSFMs match or exceed the persistence
benchmark on aggregate GDP at horizons beyond the nowcast, where classical
econometric systems typically struggle, and that they recover faster than ARIMA
under the 2008 and 2020 structural breaks. The framing of \citet{Jetwiriyanon2025}
is operational rather than theoretical, and despite the title of the paper, the
generalisation-bound claim is empirical rather than formal. The reading for our
chapter is that TSFMs interpolate well within the distribution of macroeconomic
indicators and degrade gracefully under shocks, but that the gains over the
RBNZ's own LSBoost are modest and the variance across horizons is large.

\citet{Ada2025} provide the closest thing in the literature to a TSFM in an
offline-RL loop. Their setting is a partially-observable MDP in which the
observation function $x_j(s_t) = s_t + b_j$ is offset by a per-episode unknown
$b_j$ drawn from a real-world time series, the agent does not observe $b_j$ at
test time, and the offline RL dataset was collected under the stationary case
$b_j = 0$. Their FORL framework combines two forecasting components. A
conditional diffusion model, trained on the stationary offline data, generates
$k$ candidate states from the within-episode history of observation changes and
actions. Lag-Llama, used in zero-shot mode and queried with the offsets from
previous episodes, forecasts $l$ samples of the next-episode offset. A
dimension-wise closest-match rule fuses the multi-modal diffusion candidates
with the unimodal Lag-Llama forecasts into a single state estimate that drives
a Diffusion-QL offline policy. On non-stationary maze2d and antmaze
environments under five real-world time-series offsets, FORL substantially
outperforms baselines, attaining for example $74.1$ vs $34.6$ normalised score
on maze2d-medium and $75.1$ vs $51.4$ on antmaze-umaze-diverse. The crucial
methodological observation in our setting is that the TSFM here serves only as
a passive observation-process prior. It forecasts the additive offset $b_j$ but
plays no role in modelling the action-conditional transition kernel
$s_{t+1} = f(s_t, a_t)$, which is handled by the diffusion model. The partial
identifiability argument of \citet{Ada2025}, namely that the agent cannot
disentangle $s_t$ from $b_j$ on the basis of a single observation, motivates
the diffusion-fusion step but does not itself address the deeper problem.

The composite reading of \citet{Ansari2024}, \citet{Das2024}, \citet{Rasul2023},
\citet{PerezDiaz2025}, \citet{Rahimikia2025}, \citet{Jetwiriyanon2025}, and
\citet{Ada2025} is therefore that the TSFM machinery is mature enough on the
forecasting side to serve as a passive prior in an RL pipeline, but is not yet
mature enough to serve as the action-conditional dynamics model an agent
requires. Three obstacles separate the current state from the goal. First,
pretraining is action-unconditional. Chronos, TimesFM, and Lag-Llama all model
$p(y_{T+1:T+H} \mid y_{1:T})$, possibly with exogenous covariates, but never
$p(y_{T+1:T+H} \mid y_{1:T}, a_T)$ for a decision-maker's action $a_T$. The
forecast distribution conditions on observed history, not on policy. Second,
marginal calibration does not transfer to path-dependent operational queries
(\citealp{PerezDiaz2025}, Propositions 1-3), which makes the typical CRPS-based
training objective insufficient for the threshold-crossing and run-length
quantities that economic decisions require. Third, the Lucas critique applies in
its sharpest form. A foundation forecaster trained on observational data cannot
distinguish the causal effect of a policy intervention from the response to a
confounding shock, and there is no architectural fix internal to the TSFM. The
remedies, namely identifying restrictions, instrumental variables, and
action-conditional pretraining objectives, are the subject of
Section~\ref{sec:forecasting_rl_open}.
